% !TEX root = ../algebra_02.tex

\section{Кольцо линейных операторов}

Множество всех эндоморфизмов (линейных операторов) $\End V$ образует не только линейное пространство, но и ассоциативную алгебру (кольцо с дополнительной структурой векторного пространства) относительно операций сложения и композиции.

\begin{Proposition}{Алгебра операторов и матриц}{}
	Пусть $E$ --- базис $V$, $\dim V=n$. Отображение
	\[
	\xi_E\colon \End V\to M_n(K), \qquad \xi_E(\alpha)=[\alpha]_E,
	\]
	является изоморфизмом алгебр.
	То есть оно сохраняет не только сложение и умножение на скаляр, но и композицию (которая переходит в матричное умножение):
	\[
	[\alpha+\beta]_E=[\alpha]_E+[\beta]_E, \qquad [\lambda\alpha]_E=\lambda[\alpha]_E, \qquad [\beta\alpha]_E=[\beta]_E[\alpha]_E
	\]
\end{Proposition}

\begin{Remark}{Единица кольца $\End V$}{}
	В алгебре (кольце) $\End V$ единичным элементом относительно умножения (композиции) является тождественный оператор $\operatorname{id}_V$. В соответствии с изоморфизмом $\xi_E$, ему соответствует единичная матрица $I_n$.
\end{Remark}
